% $Header: /cvsroot/latex-beamer/latex-beamer/solutions/generic-talks/generic-ornate-15min-45min.en.tex,v 1.4 2004/10/07 20:53:08 tantau Exp $

%\documentclass[trans]{beamer}
%\documentclass[draft]{beamer}
\documentclass{beamer}

\usepackage{enumerate}
\usepackage{amscd,amsmath,xspace,amssymb,amsfonts,mathrsfs,latexsym}
\usepackage{graphicx}

\usepackage{array}
%\usepackage{mbboard}
%\usepackage{multicol}
%\usepackage[pdftex]{color}

%\includeonlyframes{current}

\usepackage[all,poly,knot,tips]{xy}

\newcommand{\lra}{\longrightarrow}
\newcommand{\lla}{\longleftarrow}
\newcommand{\Lra}{\Longrightarrow}
\newcommand{\Lla}{\Longleftarrow}

%% Script Shortcuts 
\def\CA{{\mathscr A}}
\def\CB{{\mathscr B}}
\def\CC{{\mathscr C}}
\def\CD{{\mathscr D}}
\def\CE{{\mathscr E}}
\def\CF{{\mathscr F}}
\def\CG{{\mathscr G}}
\def\CH{{\mathscr H}}
\def\CI{{\mathscr I}}
\def\CJ{{\mathscr J}}
\def\CK{{\mathscr K}}
\def\CL{{\mathscr L}}
\def\CM{{\mathscr M}}
\def\CN{{\mathscr N}}
\def\CO{{\mathscr O}}
\def\CP{{\mathscr P}}
\def\CQ{{\mathscr Q}}
\def\CR{{\mathscr R}}
\def\CS{{\mathscr S}}
\def\CT{{\mathscr T}}
\def\CU{{\mathscr U}}
\def\CV{{\mathscr V}}
\def\CW{{\mathscr W}}
\def\CX{{\mathscr X}}
\def\CY{{\mathscr Y}}
\def\CZ{{\mathscr Z}}

\newcommand{\A}{\ensuremath{\mathcal{A}}\xspace}
\newcommand{\B}{\ensuremath{\mathcal{B}}\xspace}
\newcommand{\Af}{\ensuremath{\mathcal{A}_f}\xspace}
\newcommand{\C}{\ensuremath{\mathcal{C}}\xspace}
\newcommand{\D}{\ensuremath{\mathcal{D}}\xspace}
\newcommand{\F}{\ensuremath{\mathcal{F}}\xspace}
\newcommand{\I}{\ensuremath{\mathcal{I}}\xspace}
\newcommand{\K}{\ensuremath{\mathcal{K}}\xspace}
\newcommand{\Kt}{\ensuremath{\mathcal{K}_t}\xspace}
\newcommand{\Kl}{\ensuremath{\mathcal{K}_{\ell}}\xspace}
\renewcommand{\L}{\ensuremath{\mathcal{L}}\xspace}
\renewcommand{\S}{\ensuremath{\mathcal{S}}\xspace}
\newcommand{\T}{\ensuremath{\mathcal{T}}\xspace}
\newcommand{\V}{\ensuremath{\mathcal{V}}\xspace}
\newcommand{\Vf}{\ensuremath{\mathcal{V}_f}\xspace}
\newcommand{\W}{\ensuremath{\mathcal{W}}\xspace}
\newcommand{\X}{\ensuremath{\mathcal{X}}\xspace}

\newcommand{\op}{\ensuremath{^{\textrm{op}}}\xspace}
\newcommand{\ps}{\ensuremath{^{\textrm{ps}}}\xspace}
\newcommand{\lax}{\ensuremath{^{\textrm{lax}}}\xspace}
\newcommand{\oplax}{\ensuremath{^{\textrm{oplax}}}\xspace}


\newcommand{\EMK}{\textnormal{\sf EM(\K)}\xspace}
\newcommand{\EM}{\textnormal{\sf EM}\xspace}
\newcommand{\EMwK}{\ensuremath{\EM^w(\K)}\xspace}

%\newcommand{\Cat}{\textnormal{\bf Cat}\xspace}
\newcommand{\Lex}{\textnormal{\sf Lex}\xspace}
\newcommand{\Lan}{\textnormal{\sf Lan}\xspace}
\newcommand{\Hom}{\textnormal{\sf Hom}\xspace}
\renewcommand{\t}{\times}
\newcommand{\ot}{\otimes}
\newcommand{\two}{\mathbf{2}\xspace}


\newtheorem{remark}[theorem]{Remark}
\newtheorem{proposition}[theorem]{Proposition}

\beamerdefaultoverlayspecification{<+->}

%\newenvironment{\next}{\uncover<+->{}{}}

\mode<presentation>
{
%  \usetheme[compress,subsection=false]{Szeged}
  % or ...
%  \useoutertheme[subsection=false]{miniframes}
%\useoutertheme[subsection=false, footline=authortitle, frame number]{miniframes}
\useoutertheme{infolines}
\setbeamercolor{separation line}{use=structure,bg=structure.fg!50!bg}
\usecolortheme{whale}
\usecolortheme{orchid}
\usecolortheme{rose}
\setbeamercolor{titlelike}{parent=structure}


%  \setbeamercovered{transparent}
  % or whatever (possibly just delete it)
}


\usepackage[english]{babel}
% or whatever

\usepackage[latin1]{inputenc}
% or whatever

%\usepackage{times}
%\usepackage[T1]{fontenc}
% Or whatever. Note that the encoding and the font should match. If T1
% does not look nice, try deleting the line with the fontenc.


\title%[Short Paper Title] % (optional, use only with long paper titles)
{\textbf{Skew monoidality}}

%{Presentation Subtitle} % (optional)

\author{Ross Street}
\institute{Macquarie University}
%[Universities of Somewhere and Elsewhere] % (optional, but mostly needed)
%{
%  \inst{1}%
%  Department of Computer Science\\
%  University of Somewhere
%  \and
%  \inst{2}%
%  Department of Theoretical Philosophy\\
%  University of Elsewhere}
% - Use the \inst command only if there are several affiliations.
% - Keep it simple, no one is interested in your street address.

\date[July 2013]{\small Category Theory 2013, Macquarie University} 

\subject{Talks}
% This is only inserted into the PDF information catalog. Can be left
% out. 

% Delete this, if you do not want the table of contents to pop up at
% the beginning of each subsection:
\AtBeginSubsection[]
{
  \begin{frame}<beamer>
    \frametitle{Outline}
    \tableofcontents[currentsection]
  \end{frame}
}
% If you wish to uncover everything in a step-wise fashion, uncomment
% the following command: 

\beamerdefaultoverlayspecification{<.->}

\begin{document}

%  \begin{frame}
%    \titlepage
%  \end{frame}


\begin{frame}%[plain]
%  \frametitle{titlepage}

  \begin{centering}
    {\usebeamerfont{title}\usebeamercolor[fg]{title}\inserttitle

\medskip 

{\small Ross Street}

{\small Macquarie University} 

\medskip
\uncover<1->{{}}
}

\vfill

\centerline{\usebeamerfont{title}\usebeamercolor[fg]{title}\insertdate}
% CT09, Cape Town}

  \end{centering}


\end{frame}


% \begin{frame}
%   \frametitle{Outline}
%   \tableofcontents
%   % You might wish to add the option [pausesections]
% \end{frame}

%\section[Stricts and weaks]{The strict/weak dichotomy}

%\subsection{Classical universal algebra}

\begin{frame}
  \frametitle{Purpose and plan of the talk}

The goals are 
\begin{itemize}

\item to give some reasons why one might wish to weaken 
the time-honoured notion of monoidal category, and 
\pause
\item to present some theory of this weakened concept 
as developed in joint work with Steve Lack. 
\end{itemize}
\pause
I shall discuss:
 
\begin{itemize}

\item<1-> \textbf{Skew monoidal categories and skew monoidales} \\ 
\pause
\item<2-> \textbf{The connection with opmonoidal monads} \\
\pause
\item<3-> \textbf{Skew flocks} \\
\pause
\item<4-> \textbf{Flock convolution} \\
\pause
\item<5-> \textbf{Quantum categories and bialgebroids} \\
\pause
\item<6-> \textbf{Coherence for skew monoidal categories}

\end{itemize}

\end{frame}

\begin{frame}
  \frametitle{Skew monoidal categories}

A category $\CR$ is {\em right skew-monoidal} when it is equipped with 

  \begin{itemize}
\item<1-> a functor $\otimes : \CR \times \CR \lra \CR$ and object $I\in \CR$ 

\item<2-> natural families of (not necessarily invertible) morphisms 
$$\alpha : A\otimes (B\otimes C) \lra (A\otimes B)\otimes C$$ 
$$\lambda : A\lra I\otimes A$$
$$\rho : A\otimes I \lra A$$ 

\item<3-> plus five axioms as follows:

\end{itemize}

\end{frame}


\begin{frame}
  \frametitle{The pentagon and trivial unit condition}
  
\begin{equation}\label{pentagon}
\xymatrix{
& (W\otimes X)\otimes (Y\otimes Z) \ar[rd]^-{{\alpha_{W\otimes X,Y, Z}}}  & \\
W\otimes (X\otimes (Y\otimes Z)) \ar[ru]^-{{\alpha_{W, X,Y\otimes Z}}} \ar[d]_-{1_W\otimes {\alpha_{X, Y,Z}}} & & ((W\otimes X)\otimes Y)\otimes Z  \\
W\otimes ((X\otimes Y)\otimes Z) \ar[rr]_-{{\alpha_{W,X\otimes Y,Z}}} & &  (W\otimes (X\otimes Y))\otimes Z\ar[u]_-{{\alpha_{W,X,Y}} \otimes 1_Z} }
\end{equation}

\begin{equation}\label{unitunit}
\xymatrix{
I \ar[rd]_{\lambda_I}\ar[rr]^{1_I}   && I  \\
& I\otimes I \ar[ru]_{\rho_I} 
}
\end{equation}
\end{frame}

\begin{frame}
  \frametitle{Three more unit conditions}
\begin{equation}\label{leftunit}
\xymatrix{
 I\otimes (X\otimes Y)\ar[rr]^{\alpha_{I,X,Y}}   && (I\otimes X)\otimes Y  \\
& X\otimes Y \ar[lu]^{\lambda_{X\otimes Y}} \ar[ru]_{\lambda_{X} \otimes 1_Y}  &
}
\end{equation}

\begin{equation}\label{midunit}
\xymatrix{
 X\otimes (I\otimes Y)\ar[rr]^-{\alpha_{X,I,Y}} && (X\otimes I)\otimes Y \ar[d]^-{\rho_X \otimes 1_Y} \\
X\otimes Y \ar[u]^-{1_X\otimes \lambda_Y} \ar[rr]_-{1_{X\otimes Y}} && X\otimes Y}
\end{equation}

\begin{equation}\label{rightunit}
\xymatrix{
X\otimes (Y\otimes I) \ar[rr]^{\alpha_{X,Y,I}} \ar[rd]_{1_X\otimes \rho_Y}   && (X\otimes Y)\otimes I  \ar[ld]^{\rho_{X \otimes Y}} \\
& X\otimes Y   &
}
\end{equation}



\end{frame}

\begin{frame}
  \frametitle{Skew monoidales}
  A {\em right skew-monoidal} structure on an object $R$ in a monoidal bicategory $\CM$ 
consists of morphisms \hbox{$p:R\otimes R \lra R$}, $j:I\lra R$, respectively called the {\em tensor product} and {\em unit}, and 2-cells 
\begin{equation*}
\begin{aligned}
\xymatrix{
R\otimes R\otimes R \ar[d]_{1\otimes p}^(0.5){\phantom{AAAAA}}="1" \ar[rr]^{p\otimes 1}  && R\otimes R \ar[d]^{p}_(0.5){\phantom{AAAAA}}="2" \ar@{=>}"1";"2"^-{\alpha}
\\
R\otimes R \ar[rr]_-{p} && R 
}
\end{aligned}
\end{equation*}
\begin{equation*}
\begin{aligned}
\xymatrix{
R\otimes R \ar[rd]_{p}^(0.5){\phantom{aa}}="1"   && R \ar[ll]_{j\otimes 1}  \ar[ld]^{1}_(0.5){\phantom{aa}}="2" \ar@{<=}"1";"2"^-{\lambda}
\\
& R }
\end{aligned}
\qquad
\begin{aligned}
\xymatrix{
R\otimes R \ar[rd]_{p}^(0.5){\phantom{aa}}="1"   && R \ar[ll]_{1\otimes j}  \ar[ld]^{1}_(0.5){\phantom{aa}}="2" \ar@{=>}"1";"2"^-{\rho}
\\
& R }
\end{aligned}
\end{equation*}
respectively called the {\em associativity}, {\em left unit} and {\em right unit constraints}, satisfying five conditions expressible using pasting diagrams (Exercise). 
\end{frame}

\begin{frame}
  \frametitle{Example of a skew monoidal category}
  
 \begin{itemize}
\item<1-> $\mathrm{Comon}\CV$ is the monoidal category of comonoids in a nice base $\CV$.
\item<2-> Objects $A$ of $\mathrm{Comon}\CV \text{-}\mathrm{Cat} = (\mathrm{Comon}\CV) \text{-}\mathrm{Cat}$ might be considered locally non-commutative $\CV$-categories. Put $X=\mathrm{ob}A$ assumed small.
\item<3-> $\CV^X$ is the $\CV$-category of families $M=(M_x)_{x\in X}$
of objects of $\CV$. 
\item<4-> We shall define a right-skew monoidal structure on $\CV^X$ in the monoidal 2-category $\CV \text{-}\mathrm{Cat}$. 
\item<5-> The tensor product $M\ast_A N$ is defined by
$$(M\ast_A N)_x = \sum_y{M_x \otimes A(x,y) \otimes N_y} \ .$$
\item<6-> The unit object $J$ is defined by $J_x=I$ for all $x\in X$.
\end{itemize}
\end{frame}

\begin{frame}
  \frametitle{Example of a skew monoidal category (continued)}
   
$$(M\ast_A (N\ast_A L))_x = \sum_{y,z} {M_x \otimes A(x,y) \otimes N_y \otimes A(y,z)\otimes L_z}$$
$$((M\ast_A N)\ast_A L)_x = \sum_{y,z} {M_x \otimes A(x,y) \otimes N_y \otimes A(x,z)\otimes L_z} $$
 \begin{itemize}
\item<1->The morphism $\alpha : M\ast_A (N\ast_A L) \lra (M\ast_A N)\ast_A L$ uses
$$A(x,y)\otimes A(y,z) \stackrel{\delta \otimes 1}\lra A(x,y)\otimes A(x,y)\otimes A(y,z) \stackrel{1\otimes \mathrm{comp}}\lra A(x,y)\otimes A(x,z) \ .$$ 
\item<2-> 
The morphism $\lambda : N \lra J\ast_A N$ uses
$$N_x \stackrel{\mathrm{id_x}\otimes 1}\lra A(x,x)\otimes N_x \stackrel{\mathrm{in}_x}\lra \sum_y{A(x,y)\otimes N_x} \ .$$ 
\item<3-> 
The morphism $\rho : M\ast_A J\lra M$ uses $\sum_y{M_x\otimes A(x,y) \stackrel{1\otimes \varepsilon}\lra M_x}$. 
 \end{itemize}
\end{frame}

\begin{frame}
  \frametitle{Opmonoidal monads}
  An opmonoidal monad $t = (t : R\rightarrow R,\mu :tt \Rightarrow t ,\eta 1\Rightarrow t)$ 
on a right skew monoidale $R$ has 2-cells 
\begin{equation*}
\begin{aligned}
\xymatrix{
R\otimes R \ar[d]_{p}^(0.5){\phantom{AAAA}}="1" \ar[rr]^{t\otimes t}  && R\otimes R \ar[d]^{p}_(0.5){\phantom{AAAA}}="2" \ar@{=>}"1";"2"^-{\delta}
\\
R \ar[rr]_-{t} && R 
}
\qquad
\xymatrix{
I \ar[rd]_{j}^(0.5){\phantom{a}}="1" \ar[rr]^{j}  && R \ar[ld]^{t}_(0.5){\phantom{a}}="2" \ar@{<=}"1";"2"^-{\epsilon}
\\
& R 
}
\end{aligned}
\end{equation*}
making $t $, $\mu$ and $\eta$ opmonoidal.
The fusion 2-cell $\upsilon$ is 
\begin{eqnarray}\label{opmonmndfus}
\xymatrix{
R\otimes R \ar[d]_{1\otimes t}^(0.5){\phantom{AAAA}}="1" \ar[rr]^{t\otimes 1}
&& R\otimes R \ar[d]^{1\otimes t}_(0.5){\phantom{AAAAAA}}="2" \ar@{=>}"1";"2"^-{1_t \otimes \mu} \\
R\otimes R \ar[d]_{p}^(0.5){\phantom{AAAA}}="1" \ar[rr]^{t\otimes t}
&& R\otimes R \ar[d]^{p}_(0.5){\phantom{AAAAAA}}="2" \ar@{=>}"1";"2"^-{\delta} \\
R \ar[rr]_t && R }
\end{eqnarray}
\end{frame}

\begin{frame}
  \frametitle{Fusion structures}
  
Brugui\`eres-Lack-Virelizier (2011) define the opmonoidal monad to be {\em Hopf} when $\upsilon$ is invertible.

\begin{theorem}
Opmonoidal monad structures $(\mu,\eta,\upsilon,\varepsilon)$ on an endomorphism 
$t$ of a right skew monoidale $R$ are in bijection,
via the definition \eqref{opmonmndfus} of $\upsilon$, with triples $(\upsilon,\eta, \varepsilon)$ 
of 2-cells satisfying five specific axioms.
\end{theorem}
Here are the five axioms.
\begin{equation*}\label{fusionax4&5}
\begin{aligned}
\xymatrix{
p(1\otimes t) \ar[rd]_{p(\eta \otimes 1)}\ar[rr]^{\eta p(1\otimes t)}   && tp(1\otimes t) \ar[ld]^{\upsilon} \\
& p(t\otimes t)  &
}
\end{aligned}
\qquad
\begin{aligned}
\xymatrix{
j\ar[rd]_{1_j}\ar[rr]^{\eta j}   && tj \ar[ld]^{\varepsilon} \\
& j & 
}
\end{aligned}
\end{equation*}
\end{frame}

\begin{frame}
\begin{equation*}\label{fusionax2}
\begin{aligned}
\xymatrix{
& p(t \otimes t)(1\otimes j) \ar[rd]^-{\phantom{A}\cong}  & \\
tp(1\otimes t)(1\otimes t) \ar[ru]^-{\upsilon (1\otimes j)\phantom{A}} \ar[d]_-{tp(1\otimes \varepsilon)} & & p(1\otimes t)(1\otimes j)t \ar[d]^-{p(1\otimes \varepsilon)t} \\
tp(1\otimes j) \ar[r]_-{\rho} & t & p(1\otimes j)t \ar[l]^-{\rho}    
}
\end{aligned}
\end{equation*}

\begin{equation*}\label{fusionax3}
\begin{aligned}
\xymatrix{
tp(1\otimes t)(j\otimes 1) \ar[rr]^{\upsilon (j\otimes 1)} & &p(t\otimes t)(j\otimes 1) \ar[d]^{\cong} \\
tp(j\otimes 1) \ar[u]^{tp(1\otimes \eta)(j\otimes 1)} & & p(tj\otimes 1)t \ar[d]^{p(\varepsilon \otimes 1)t} \\
t \ar[rr]_{\lambda t}  \ar[u]^{t\lambda} & & p(j\otimes 1)t 
}
\end{aligned}
\end{equation*}
\end{frame}
\begin{frame}
\begin{equation*}\label{fusionax1}
\begin{aligned}
\xymatrix{
tp(1\otimes t)(1\otimes p)(1\otimes 1\otimes t) \ar[d]_-{tp(1\otimes \upsilon)} \ar[rr]^-{\upsilon(1\otimes p)(1\otimes 1\otimes t)} & & p(t\otimes t)(1\otimes p)(1\otimes 1\otimes t) \ar[d]^-{p(t\otimes 1)(1\otimes \upsilon)}\\
tp(1\otimes p)(1\otimes t\otimes t) \ar[d]_-{\alpha}  & & p(t\otimes 1)(1\otimes p)(1\otimes t\otimes t) \ar[d]^-{\alpha}  \\
tp(1\otimes t)(p\otimes 1)(1\otimes t\otimes 1) \ar[d]_-{\upsilon (p\otimes 1)(1\otimes t\otimes 1)} & & p(p\otimes 1)(t\otimes t\otimes t)  \\
p(t\otimes t)(p\otimes 1)(1\otimes t\otimes1) \ar[rr]_{ p(1\otimes t)( \upsilon \otimes 1)} & &  p(1\otimes t)(p\otimes 1)(t\otimes t\otimes 1)  \ar[u]_-{\cong} \\
}
\end{aligned}
\end{equation*}

\end{frame}

\begin{frame}
  \frametitle{Skew flocks}
A {\itshape skew flock} in a monoidal bicategory $\CM$ is an object $R$ equipped
with a right bidual $R^{\circ }$, having unit and counit $n:I\longrightarrow
R^{\circ }\otimes R$ and $e:R\otimes R^{\circ }\longrightarrow I$,
a morphism
\begin{equation*}
q:R\otimes R^{\circ }\otimes R\longrightarrow R%
\end{equation*}
and 2-cells
\begin{gather}
\begin{split}
\xymatrix{
R\otimes R^{\circ} \otimes R\otimes R^{\circ} \otimes R \ar[d]_-{1\otimes 1\otimes q} \ar[rr]^-{q\otimes 1 \otimes 1}_-{\phantom{A}}="1"  &&  R\otimes R^{\circ} \otimes R \ar[d]^-{q} \\
R\otimes R^{\circ} \otimes R \ar[rr]_-{q}^-{\phantom{A}}="2" \ar@{<=}"1";"2"_{\phi} && R
} 
\end{split}
\label{phi}
\end{gather}
\end{frame}
\begin{frame}
\begin{gather}
\begin{split}
\xymatrix{
\ar@{}[rrrr]_{\phantom{A}}="1" && R\otimes R^{\circ} \otimes R \ar[rrd]^{q} && \\
R \ar[rru]^{1\otimes n} \ar[rrrr]_-{1}^{\phantom{a}}="2" \ar@{=>}"1";"2"^{\xi}  &&  &&  R 
}
\end{split}
\label{xi}
\end{gather}
\begin{equation}
\begin{split}
\xymatrix{
R \otimes R^{\circ} \otimes R \ar@/^/@<+1ex>[rrrr]^-{e\otimes 1}_-{}="1" \ar@/_/@<-1ex>[rrrr]_-{q}^-{}="2" \ar@{=>}"1";"2"^{\zeta} &&&& R 
}
\end{split}
\label{zeta}
\end{equation}
\end{frame}
\begin{frame}
\noindent satisfying the following five conditions (where $1_{n}=\overbrace{1\otimes . . . \otimes 1}^{n}$):
\begin{gather*}
\begin{split}
\scalebox{0.7489}{
\xymatrix{
\ar@{}[rr]_-{\phantom{a}}="1" & q(1_{2}\otimes q)( q\otimes 1_{4}) \cong q(q\otimes 1_{2})( 1_{4}\otimes q) \ar@{<-}[rd]^-{\phantom{AAA}\phi (1_{4}\otimes q)} & \\
(q( q\otimes 1_{2})(q\otimes 1_{4}) \ar@{<-}[ru]^-{\phi (q\otimes 1_{4})\phantom{AAA}} \ar@{<-}[d]_-{q(\phi\otimes 1_{2})\phantom{A}} && q( 1_{2}\otimes q)( 1_{4}\otimes q)) \\
q( q\otimes 1_{2})(1_{2}\otimes q\otimes 1_{2}) \ar@{<-}[rr]_-{\phi(1_{2}\otimes q\otimes 1_{2})}="2" && q( 1_{2}\otimes q)( 1_{2}\otimes q\otimes 1_{2})  \ar@{<-}[u]_{\phantom{A}q( 1_{2}\otimes \phi)}
\ar@{}"1";"2"|*=0-{}
}
}
\end{split}
\end{gather*}
\end{frame}
\begin{frame}
\begin{gather*}
\begin{split}
\xymatrix{
1 \ar@{ ->}@/_3ex/@<-0.7ex>[rrrd]_-{1}="2"  \ar@{}[r]|-{\small{\cong}}  & (e\otimes 1)( 1\otimes n) \ar@{->}[r]^-{\zeta (1\otimes n)}="1" & q( 1\otimes n) \ar@{->}[rd]^-{\xi} & \\
 &&& 1 \ar@{}"1";"2"|*=0-{}
}
\end{split}
\end{gather*}
\begin{equation*}
\xymatrix{
(e\otimes1)(1_2\otimes q)\ar[rr]^-{\zeta (1_2\otimes q)} \ar[d]_-{\cong} && q(1_2 \otimes q) \ar[d]^-{\phi} \\
q(e\otimes 1_3) \ar[rr]_-{q(\zeta \otimes 1_2)} && q(q\otimes 1_2)}
\end{equation*}
\begin{equation*}
\xymatrix{
q(1_2\otimes q)(1_3\otimes n) \ar[rr]^-{\phi (1_3\otimes n)} \ar[d]_-{q(1_2\otimes \xi)} && q(q\otimes 1_2)(1_3\otimes n) \ar[d]^-{\cong} \\
q  && q(1\otimes n)q \ar[ll]^-{\xi q}}
\end{equation*}
\end{frame}
\begin{frame}
\begin{gather*}
\begin{split}
\xymatrix{
 q(1_{2}\otimes q)( 1_{1}\otimes n\otimes 1_{2}) \ar@{->}[r]^{\phi (1_{1}\otimes n\otimes 1_{2})} & q( q\otimes 1_{2})( 1_{1}\otimes n\otimes 1_{2}) \ar@{->}[d]^-{\phantom{A}q(\xi \otimes 1_{2})}="2" \\
q( 1_{2}\otimes e\otimes 1_{1})(1_{1}\otimes n\otimes 1_{2}) \cong q 1_{3} \cong q  \ar@{->}[u]^-{q( 1_{2}\otimes \zeta)( 1_{1}\otimes n\otimes 1_{2})\phantom{A}}="1" \ar@{->}[r]_-{1} & q\cong q 1_{3}
\ar@{}"1";"2"|*=0-{\phantom{....}}
}
\end{split}
\end{gather*}
\bigskip
A {\em flock} as per Booker-Street (2011) is a skew flock with $\phi$ invertible.
\end{frame}
\begin{frame}
  \frametitle{Example of a skew flock}
  \begin{itemize}
\item Take any tensored $\CV$-category $\CA$:
$$\pi : \CA (X\star B,C) \cong [X, \CA (B,C)] \ .$$
\pause
\item Put $\mathrm{ev} =\pi^{-1}(1_{\CA (B,C)}): \CA (B,C)\star B \lra C$ and $\mathrm{id}=\pi (I\star B \cong B) : I \lra \CA(B,B)$.
\pause
\item Let $Q : \CA \otimes \CA^{\mathrm{op}} \otimes \CA \lra \CA$ be the functor
$$Q(A,B,C) = \CA (B,A)\star C \ .$$
\pause
\item 
$\pi((\CA (B,A) \otimes \CA (D,C))\star D \cong \CA (B,A) \star (\CA (D,C)\star D)  \stackrel{1\otimes \mathrm{ev} }\lra \CA (B,A)\star C)$ 
$$= (\CA (B,A)\otimes \CA (D,C) \lra \CA (D,\CA (B,A)\star C))$$
to which we apply $-\star E$ to obtain a natural family
$$\phi : Q(A,B,Q(C,D,E)) \lra Q(Q(A,B,C),D,E) \ .$$ 
 \end{itemize}
\end{frame}
\begin{frame}
 \begin{itemize}
\item The evaluation $\mathrm{ev} : \CA (A,B)\star A \lra A$ gives a natural family 
$$\xi : Q(B,A,A) \lra B \ .$$
\pause
\item $\mathrm{id}\otimes 1 : A \lra \CA (B,B)\star A$ gives a natural family
$$\zeta : A \lra Q(B,B,A) \ .$$
This defines a skew flock $(\CA,Q_*, \phi, \xi, \zeta)$ in the putative monoidal bicategory $\CV\text{-}\mathrm{Mod}$.
 \end{itemize}
\end{frame}
\begin{frame}
 \frametitle{Skew flocks and opmonoidal monads}
 \begin{itemize}
 \item A biduality $R\dashv _{b}R^{\circ }$ generates a monoidal structure on $R^{\circ }\otimes R$. 
 \pause
 \item Given a skew flock $R$, we have the mate
\begin{equation*}
\hat{q}:R^{\circ }\otimes R\longrightarrow R^{\circ }\otimes R
\end{equation*}
of $q:R\otimes R^{\circ }\otimes R\longrightarrow R$ under the biduality $R\dashv _{b}R^{\circ }$.
\end{itemize}
\pause
\begin{proposition}\label{opmnd=skflk} 
There is a bijection between skew flock structures on $(R,q)$ 
and opmonoidal monad structures on $(R^{\circ }\otimes R, \hat{q})$.  
\end{proposition}
\end{frame}
\begin{frame}
 \frametitle{Skew flocks and skew monoidales}
 \begin{itemize}
 \item Let $(R,q)$ be a skew flock in the monoidal bicategory $\CM$.
 \pause
 \item Let $j : I \lra R$ have a right adjoint $j^*$ with counit $\varepsilon_j : j\circ j^* \Lra 1_R$.
 \pause
\item Put $h = j^{* \circ} : I \lra R^{\circ}$. 
Then $h$ also has a right adjoint $h^* = j^{\circ}$.
Moreover, $j^* = e\circ (1\otimes h)$,  $h \cong (1\otimes j^*)\circ n$
and 
$j \circ j^* \cong (j^* \otimes 1_R)\circ (1_R\otimes j) \cong  j^* \otimes j $.
So $(h\otimes 1)\circ j \cong j\circ j^* \circ n$.
\pause
\item Define
\begin{eqnarray*}
p = (R\otimes R \stackrel{1\otimes h \otimes 1}\lra R\otimes R^{\circ} \otimes R \stackrel{q}\lra R) \ .
\end{eqnarray*}
\pause
\item Obtain $\alpha : p(1\otimes p) \Lra p(p\otimes 1)$ from
$\phi (1\otimes h\otimes 1\otimes h\otimes 1) :$
\begin{eqnarray*}
q(1\otimes 1\otimes q)(1\otimes h\otimes 1\otimes h\otimes 1)
\Lra 
q(q\otimes 1\otimes 1)(1\otimes h\otimes 1\otimes h\otimes 1)
\end{eqnarray*}
and the canonical isomorphisms
$$p(1\otimes p) \cong q(1\otimes 1\otimes q)(1\otimes h\otimes 1\otimes h\otimes 1)$$ and  
$$q(q\otimes 1\otimes 1)(1\otimes h\otimes 1\otimes h\otimes 1) \cong p(p\otimes 1) \ .$$
 \end{itemize}
 \end{frame}
 \begin{frame}
 \begin{itemize}
\item Next $\lambda : 1_R \Lra p\circ (j\otimes 1_R)$ has mate
$j^*\otimes 1 \Lra p$ under the adjunction $j\otimes 1\dashv j^*\otimes 1$ equal to
$$\zeta \circ (1\otimes h\otimes 1) : (e\otimes 1)\circ (1\otimes h\otimes 1) \Lra q\circ (1\otimes h\otimes 1) \ .$$ 
\pause
\item Finally $\rho : p\circ (1_R\otimes j) \Lra 1_R$ is the composite 
$$q\circ (1\otimes h\otimes 1)\circ (1\otimes j)\cong q\circ (1\otimes (j\circ j^*))\circ (1\otimes n) \Lra q\circ (1\otimes n) \stackrel{\xi}\Lra 1 $$
where the unlabelled 2-cell is $q\circ (1\otimes \varepsilon_j)\circ (1\otimes n)$.
\end{itemize}
\pause
\begin{proposition}\label{smfromsf}
For any skew flock $R$ and any left adjoint morphism $j : I \lra R$ in $\CM$, with 
$p$, $\alpha$, $\lambda$ and $\rho$ as above, $(R, p, j)$ is a right skew monoidale in $\CM$.  
\end{proposition}
 \end{frame}
 \begin{frame}
 \begin{Definition} A {\em right skew $\CV$-proflock} is a right skew flock $\CA$ in $\CV\text{-}\mathrm{Mod}$.
\end{Definition}
 
A {\em right skew} $\CV$-{\em flock} is a right skew $\CV$-proflock $\CA$ in $\CV\text{-}\mathrm{Mod}$ for which the module\ \
\begin{equation}
Q:\CA\otimes \CA^{\mathrm{op}}\otimes \CA\longrightarrow
\CA
\end{equation}
is (representable as) a $\CV$-functor. 
More explicitly, we have a $\CV$-category $\CA$ and a
$\CV$-functor $Q:\CA\otimes \CA^{\mathrm{op}}\otimes
\CA\longrightarrow \CA$ equipped with $\CV$-natural
transformations
\begin{gather*}
\phi :Q( A,B,Q(C,D,E)) \longrightarrow Q(Q(A,B,C),D,E)
\end{gather*}
\begin{gather*}
\xi _{A}^{B}:Q( A,B,B) \longrightarrow A
\end{gather*}
\begin{equation*}
\zeta _{B}^{A}:B\longrightarrow Q( A,A,B) \ ,
\end{equation*}
\end{frame}
\begin{frame}
such that the following five conditions hold:
\begin{gather*}
\begin{split}
\xymatrix{
Q(Q(A,B,C),D,Q(E,F,G)) \ar[d]_{\phi}="1" & Q(A,B,Q(C,D,Q(E,F,G))) \ar[l]_{\phi} \ar[dd]^{Q(1,1,\phi)}\\
Q(Q(Q(A,B,C),D,E),F,G)  &  \\
Q(Q(A,B,Q(C,D,E)),F,G)\ar[u]^{Q(\phi,1,1)}  & Q(A,B,Q(Q(C,D,E),F,G) \ar[l]^{\phi}
} 
\end{split}\\ 
\begin{split}
\xymatrix{
Q(A,B,Q(B,B,C)) \ar[rr]^-{\phi} && Q(Q(A,B,B),B,C) \ar[d]^{Q(\xi,1,1)}="2" \\
Q(A,B,C) \ar[u]^{Q(1,1,\zeta)}="1" \ar[rr]_{1} && Q(A,B,C)
} 
\end{split}
\end{gather*}
\end{frame}
\begin{frame}
\begin{equation*}
\xymatrix{
Q(A,A,Q(B,C,D)) \ar[rr]^{\phi}   && Q(Q(A,A,B),C,D)  \\
& Q(B,C,D) \ar[lu]^{\zeta} \ar[ru]_{Q(\zeta,1,1)}  &
}
\end{equation*}
\begin{equation*}
\xymatrix{
Q(A,B,Q(C,D,D)) \ar[rd]_{Q(1,1,\xi)}\ar[rr]^{\phi}   && Q(Q(A,B,C),D,D) \ar[ld]^{\xi} \\
& Q(A,B,C)  &
}
\end{equation*}
\begin{gather*}
\begin{split}
\xymatrix{
&& Q(A,A,A)\ar[rrd]^{\xi} && \\
A \ar[rru]^{\zeta} \ar[rrrr]_-{1} &&  &&  A \ . \ \square 
} 
\end{split}
\end{gather*}
 \end{frame}
 \begin{frame}
  \frametitle{Flock convolution}
 \begin{itemize}
\item $\CV\text{-}\mathrm{Cocts}$ is the 2-category of cocomplete $\CV$-categories
and colimit-preserving (= continuous) $\CV$-functors made monoidal via $\CX \boxtimes \CY$:
$\CV\text{-}\mathrm{Cocts}(\CX \boxtimes \CY,\CZ) \simeq \CV\text{-}\mathrm{Sepcocts}(\CX \otimes \CY,\CZ)$.
\pause
\item $\CV\text{-}\mathrm{Mod}(\CI,-) : \CV\text{-}\mathrm{Mod} \lra \CV\text{-}\mathrm{Cocts} $
is a strong monoidal pseudofunctor taking $\CA$ to $[\CA^{\mathrm{op}},\CV]$.
\pause
\item So it takes right skew $\CV$-proflocks $\CA$ to right skew flocks $[\CA^{\mathrm{op}},\CV]$
in $\CV\text{-}\mathrm{Cocts}$. 
\pause
\item The basic formula is
\begin{equation}
Q(M,N,L)A=\int^{B,C,D}{Q(A;B,C,D)\otimes MB\otimes NC\otimes LD}
\end{equation}
yielding a $\CV$-functor
\begin{equation}
Q:[\CA^{\mathrm{op}},\CV]\otimes [\CA,\CV]\otimes [\CA^{\mathrm{op}},\CV]\lra [\CA^{\mathrm{op}},\CV]
\end{equation}
which is colimit-preserving in each variable separately.
\end{itemize}
\end{frame}
\begin{frame}
  \frametitle{Many closed skew-monoidal categories}
According to our previous Proposition, this gives a multitude of examples of closed skew-monoidal $\CV$-categories 
$[\CA^{\mathrm{op}},\CV]$, 
one for each choice of object $K$ in $\CA$ (or its Cauchy completion); namely,
we take the tensor product
$$M\ast N = Q(M,\CA (K,-),N)$$
with skew unit $\CA (-,K)$.
 \end{frame}
 \begin{frame}
  \frametitle{Abstract quantum categories and bialgebroids}
\begin{itemize}
\item Under a different name, bialgebroids occurred in the work of Takeuchi (1977). 
They arose again in the quantum group literature.
They were related to opmonoidal monads by Szlach\'anyi (2003) and put into the context
of monoidal bicategories by Day and the author (2004). 
\pause
\item A {\em bialgebroid} in a monoidal bicategory $\CM$ is a biduality 
$R \dashv R^{\circ}$ equipped with an opmonoidal monad on the generated monoidale 
$R^{\mathrm{e}}=R^{\circ} \otimes R$ and a distinguished map $j:I\lra R$. 
\pause 
\item Our earlier Propositions show this is the same as a skew flock $R$ in $\CM$ with a distinguished map $j:I\lra R$, and this gives rise to a skew monoidale $R$. 
\pause
\item Lack-Street (2012) shows that the skew flock can be recaptured 
from the skew monoidale when the right adjoint $j^*$ of $j$ has a persistent monadicity property.
In the classical examples of $\CM$, there is a canonical choice of $j : I \lra R$.  
\pause
\item A bialgebroid is {\em Hopf} when the fusion morphism of the opmonoidal monad is invertible;
that is, when the skew flock is a flock.  
\end{itemize}

\begin{example} 
Let $\mathrm{Mod}\CV$ denote the monoidal full sub-bicategory of $\CV\text{-}\mathrm{Mod}$ 
consisting of the one-object $\CV$-categories. So the objects $R$ of $\mathrm{Mod}\CV$ are 
monoids in $\CV$. There is the monoid morphism $\eta : I \lra R$ which is the unit of $R$
and so provides a distinguished map $\eta_* : I \lra R$ in $\mathrm{Mod}\CV$. 
The classical bialgebroids are bialgebroids $R$ in the bicategory $\mathrm{Mod}\CV$ with $j = \eta_*$ where $\CV$ is the monoidal category of modules over a commutative ring $k$. 
Hopf bialgebroids are sometimes called ``quantum groupoids''.   $\square$   
\end{example}

\begin{example} 
Suppose the tensor product in $\CV$ preserves coreflexive equalizers separately in each variable.
Let $\mathrm{Comod}\CV$ denote the monoidal bicategory whose objects $C$ are 
comonoids in $\CV$ and whose morphisms are two-sided comodules; see \cite{QCat}. 
There is the comonoid morphism $\varepsilon : C \lra I$ which is the counit of $C$
and so provides a distinguished map $\varepsilon^* : I \lra C$ in $\mathrm{Comod}\CV$. 
Quantum categories are bialgebroids $C$ in the bicategory $(\mathrm{Comod}\CV)^{\mathrm{co}}$ with $j = \varepsilon^*$ where $\CV$ is the monoidal category of modules over a division ring $k$.
Hopf quantum categories are {\em quantum groupoids} as per \cite{ChLaSt}.   $\square$
\end{example}

The following was proved by Lack and the author in \cite{EckHil}. 

\begin{theorem}
For a cartesian bicategory $\CM$ in the sense of \cite{CBII}, 
skew monoidales $C$ in $\CM$ with unit the unique map $I\to C$, 
amount to monads on $C$ in $\CM$.
\end{theorem}

The bicategory $\mathrm{Span}$ of spans of sets can be generalized to 
a bicategory $\mathrm{Span}(\CE)$ of spans in a finitely complete 
category $\CE$; taking $\CE$ to be the category of sets and functions, 
we recover $\mathrm{Span}$ itself.
The bicategory $\mathrm{Span}(\CE)$ is cartesian and so
in there quantum categories are just monads, 
and, as noted by B\'enabou in \cite{Ben1967}, 
these are just categories internal to $\CE$.
\end{frame}

 \begin{frame}
  \frametitle{Bialgebroids}
\begin{example} 
\begin{itemize}
\item Let $\mathrm{Mod}\CV$ denote the monoidal full sub-bicategory of $\CV\text{-}\mathrm{Mod}$ 
consisting of the one-object $\CV$-categories. So the objects $R$ of $\mathrm{Mod}\CV$ are 
monoids in $\CV$. 
\pause
\item There is the monoid morphism $\eta : I \lra R$ which is the unit of $R$
and so provides a distinguished map $\eta_* : I \lra R$ in $\mathrm{Mod}\CV$. 
\pause
\item The classical bialgebroids are bialgebroids $R$ in the bicategory $\mathrm{Mod}\CV$ with $j = \eta_*$ where $\CV$ is the monoidal category of modules over a commutative ring $k$. 
\pause
\item Hopf bialgebroids are sometimes called ``quantum groupoids''.  
\end{itemize}  
\end{example}
\end{frame}
 \begin{frame}
  \frametitle{Quantum categories}
\begin{example} 
\begin{itemize}
\item Suppose the tensor product in $\CV$ preserves coreflexive equalizers separately in each variable.
\pause
\item Let $\mathrm{Comod}\CV$ denote the monoidal bicategory whose objects $C$ are 
comonoids in $\CV$ and whose morphisms are two-sided comodules. 
\pause
\item There is the comonoid morphism $\varepsilon : C \lra I$ which is the counit of $C$
and so provides a distinguished map $\varepsilon^* : I \lra C$ in $\mathrm{Comod}\CV$. 
\pause
\item Quantum categories are bialgebroids $C$ in the bicategory $(\mathrm{Comod}\CV)^{\mathrm{co}}$ with $j = \varepsilon^*$ where $\CV$ is the monoidal category of modules over a division ring $k$.
\pause
\item Hopf quantum categories are {\em quantum groupoids} as per Chikhladze-Lack-Street (2010).
\end{itemize} 
\end{example}
\end{frame}

\begin{frame}
\frametitle{The Message}

\Large{
\[
\frac{\mathbf{Quantum \ \ groupoids}}{\mathbf{Flocks}}=\frac{\mathbf{Groups}}{\mathbf{Herds}}
\]}
\end{frame}

 \begin{frame}
  \frametitle{The cartesian case}

\begin{theorem} [Lack-Street]
For a cartesian bicategory $\CM$ in the sense of Carboni-Kelly-Walters-Wood (2008), 
skew monoidales $C$ in $\CM$ with unit the unique left adjoint morphism $I\to C$, 
amount to monads on $C$ in $\CM$.
\end{theorem}
\bigskip
The bicategory $\mathrm{Span}$ of spans of sets can be generalized to 
a bicategory $\mathrm{Span}(\CE)$ of spans in a finitely complete 
category $\CE$; taking $\CE$ to be the category of sets and functions, 
we recover $\mathrm{Span}$ itself.
The bicategory $\mathrm{Span}(\CE)$ is cartesian and so
in there quantum categories are just monads, 
and, as noted by B\'enabou (1967), 
these are just categories internal to $\CE$.
\end{frame}

 \begin{frame}
  \frametitle{Coherence for skew monoidal categories}
  \begin{itemize}
\item Left skew-monoidal categories make a clubbable situation in the sense of Kelly (1972). 
\pause
\item Write $\mathbf{Fsk}$ for the free left skew-monoidal category on a single generating object $N$. 
We have a concrete model for this.
\pause
\item Let $\mathbf{\Delta}_{\bot}$ denote the category whose objects are the finite ordinals 
$\mathbf{n} = \{0,1,\dots , n-1\}$
with first element (so that $n>0$), 
and whose morphisms are first-element-and-order-preserving functions.
\pause
\item $\mathbf{\Delta}_{\bot}$ is Hopf left skew monoidal with ordinal sum as
tensor product, with unit $\mathbf{1}$, with left unit constraint $\lambda : \mathbf{1+n}\lra \mathbf{n}$
the surjection identifying $0$ and $1$, and with right unit constraint $\rho : \mathbf{m}\lra \mathbf{m+1}$ the injection whose image does not contain $m+1$.   
\end{itemize}
\pause
\begin{theorem}[Lack-Street]
The unique strict monoidal functor $\mathbf{Fsk} \lra \mathbf{\Delta}_{\bot}$, which takes the 
generating object $N$ of $\mathbf{Fsk}$ to the object $\mathbf{1}$ of 
$\mathbf{\Delta}_{\bot}$, is faithful.  
\end{theorem}
\end{frame}
\end{document}
